% Pursuing Stacks (cotes 134-2 to 134-8): archive cote 134-5, site page 105 (PDF page 106), typescript page 292.
% typescript. Draft compared: ps5.tex:52-116.
% AI-transcribed by Claude Opus 5.5 (claude-opus-5-5, adaptive thinking, effort high), Message Batches API, batch ps2 (msgbatch_01NVwr1Yuo8qJf5rykeGBXPK), 2026-11-12; not yet checked by eye.
\scanpage{134-5/105}
\mspage{292}

what his assumptions on the indexing categories are.  Thomason seems to
believe that the closed model structure of (Cat) is essential for being
able to take homotopy limits  -  whereas it is clear a priori to me that
the notion \xout{should} depends only on the notion of weak equivalence.  Indeed,
he seems to consider the possibility of taking homotopy limits in (Cat)
as the main application of his theorem,  and in order to apply Anderson's
results,  believes it is necessary to be able to give concrete caracteri-%
zations for ``fibrations'' and ``cofibrations'' of his closed model structure.
Now it turns out that the case he is interested in (for proving ``Lichten-%
baum's conjecture'') is \xout{the case of} a typical case of \uline{direct} homotopy
limits,  namely ``\uline{integration}''  -  which can be described directly in (Cat)
in such an amazingly simple way (as sketched in section 69, p.198-199).
Thus I feel for this application,  the closed model structure is wholly
irrelevant.  As for cointegration,  I do not expect that there is a com-%
parably simple construction of this operation  within the modelizer (Cat),
but presumably there is when taking $\infty$-Gr-stacks as models  (as suggested
by the ``geometric'' approach to cohomology invariants,  via stacks, \xout{precisely}
where the operation of ``direct image''  namely cohomology precisely,  is
the obvious one,  whereas inverse images are more delicate to define, by
an adjunction property with respect to direct images ...).  When working
in (Cat),  cointegration of homotopy types should be not more nor less
involved than in  any closed model category say,  and involve intensive
recourse to ``fibrations'' (in the sense of the closed model structure,  or
more intrinsically,  in the sense that base change by these should preser-%
ve weak equivalence).  Now  the latter have become quite familiar \inserted[typed]{to me}\xout{\unsure{to me}}
during my long scratchwork on cohomology properties of maps in (Cat),
and I'll \xout{repeat} have to try put it down nicely in part V of the notes.

88.     The present ``part IV''  on asphericity structures (and their rela-%
tions to contractibility structures)  turns out a lot longer than I had
anticipated,  and the end is not yet quite in sight !  Therefore,  before
pursuing,  I would like to make a review of the questions \xout{which} along
these lines which seem to require elucidation,  and then decide which I'm
going to deal with,  before going over to part V .

1)   Whereas the relevant notions of ``morphisms''  and ``bimorphisms''
for contractibility structures seem to me well understood,  there remains

\margin{105}
